% Before start, check the file README.txt
\documentclass[10pt,a4paper,twocolumn,english]{article}
\usepackage{sbsr}
\usepackage{float}

%%%% This file presents the global settings for the document in LaTeX %%%%
%%%% (MORE INFORMATION IN README.txt) %%%%
%%%  The paper information must be declared below as well as its authors and affiliations %%%%

% analisis exploratorio de un conjunto de datos vis-nir para segmentacion y clasificacion de cobertura y uso del suelo mediante modelos de inteligencia artificial no supervisada 
%agregar kmeans al titulo
\title{Cross-Domain Spectral Compatibility and Unsupervised Land-Cover Classification of VIS-NIR UAV Imagery in Colombian Landscapes}

\author{Juan Angel Gamez Diaz$^1$}

\address{
    Sebastián Carmona Estrada, Tutor, EAFIT University, School of Applied Sciences and Engineering, Medellín, Colombia – scarmonae@eafit.edu.co; 
    Olga Lucía Quintero Montoya, Co-tutor, EAFIT University, School of Applied Sciences and Engineering, Medellín, Colombia –oquinte1@eafit.edu.co
}

%%%% Beginning of the document %%%%

\begin{document}


% create the title %%%% (DO NOT CHANGE) %%%%
\twocolumn[
\begin{@twocolumnfalse}
\maketitle
\end{@twocolumnfalse}
]

%en que aporta esto a los resultados.
%%%% Add each section file %%%%
%%%% (ALTER ONLY IF SECTIONS ARE INTRODUCED/REMOVED) %%%%

\begin{abstract}

    Diabetic retinopathy screening is limited by costly fundus imaging equipment.
    Super-resolution offers a computational path to recover image quality from low-cost sensors.
    This work presents an SRGAN framework for retinal fundus enhancement under realistic degradation.
    Nine degradation scenarios are defined: three downsampling methods (bicubic, bilinear, Lanczos), two noise models (Gaussian, Poisson), and three scale factors ($\times$2, $\times$4, $\times$8).
    The pipeline follows a camera-motivated order: blur $\rightarrow$ downsample $\rightarrow$ noise.
    The generator uses a 16-block SRResNet with PixelShuffle upsampling.
    Training proceeds in two phases: L1 pre-training followed by full adversarial optimization with perceptual (VGG-19), adversarial, and pixel-wise losses.
    A VGG16 classifier is trained on super-resolved images for DR severity grading.
    Best results reach 37.43\,dB PSNR, 0.944 SSIM, and 0.043 LPIPS at $\times$2 scale.
    Scenario B (bilinear) outperforms bicubic and Lanczos across all metrics.
    The framework provides a reproducible benchmark for retinal SR under controlled degradation.

\end{abstract}


\section{Introduction}

Diabetic retinopathy (DR) is a leading cause of preventable blindness~\cite{gulshan2016}.
Early screening reduces severe vision loss risk, but high-resolution fundus cameras remain unaffordable in low-resource settings~\cite{haque2017}.
Low-cost imaging alternatives produce degraded images with limited diagnostic value~\cite{willemink2020}.

Deep learning-based super-resolution (SR) can recover high-frequency details from low-resolution inputs~\cite{zhang2020}.
Generative adversarial networks (GANs)~\cite{goodfellow2014} achieve photorealistic upsampling through adversarial training.
The SRGAN architecture~\cite{ledig2017} combines perceptual, adversarial, and pixel-wise losses for high-quality reconstruction.
ESRGAN~\cite{wang2018} further improves visual quality through residual-in-residual dense blocks.

Despite these advances, three gaps persist.
First, public fundus datasets are collected with high-end cameras and do not represent low-cost acquisition~\cite{willemink2020}.
Second, existing retinal SR studies evaluate single degradation models, preventing systematic comparison.
Third, the impact of SR on downstream DR classification remains underexplored within unified frameworks.

This work addresses these gaps with four contributions.
First, a 9-scenario experimental design spanning three interpolation methods, two noise types, and three scale factors.
Second, a complete SRGAN implementation with two-phase training and perceptual loss.
Third, quantitative evaluation with PSNR, SSIM, and LPIPS across all scenarios.
Fourth, a VGG16 DR classifier on super-resolved images for end-to-end assessment.

The paper is organized as follows.
Section~\ref{sec:methods} describes the dataset, degradation pipeline, and architecture.
Section~\ref{sec:results} presents quantitative and qualitative results.
Section~\ref{sec:discussion} interprets the findings.
Section~\ref{sec:conclusions} concludes and outlines future work.


\section{Materials and Methods}
\label{sec:methods}

\subsection{Dataset}

The dataset combines APTOS 2019 and Messidor-2 retinal fundus collections.
All images are preprocessed into a unified high-resolution reference set (\textbf{M3}) of 3,578 RGB photographs at 512$\times$512 pixels.
The set is split into training (70\%), validation (15\%), and test (15\%) subsets.
DR severity spans five clinical grades: none (0), mild (1), moderate (2), severe (3), and proliferative (4)~\cite{gulshan2016}.
Degradations are applied on-the-fly during training, eliminating disk storage of multiple copies and ensuring reproducibility through fixed seeds.

\subsection{Degradation Pipeline}

The pipeline simulates low-cost camera image formation.
It follows the physically motivated order~\cite{ledig2017}:

\begin{equation}
\mathbf{I}_{\text{LR}} = \left( \mathbf{I}_{\text{HR}} * k_\sigma \right) \downarrow_s + n
\label{eq:degradation}
\end{equation}

where $k_\sigma$ is a Gaussian blur kernel, $\downarrow_s$ is downsampling by scale $s$, and $n$ is noise.
The order models: lens optics (blur) $\rightarrow$ sensor discretization (downsample) $\rightarrow$ electronics (noise).

Three scenarios are defined (Table~\ref{tab:experiments}).
Scenario A (bicubic, $\sigma=1.0$, $\sigma_n=5$) is the academic baseline.
Scenario B (bilinear, $\sigma=1.5$, $\sigma_n=10$) simulates mid-range portable cameras.
Scenario C (Lanczos, $\sigma=2.0$, Poisson) models extreme low-light acquisition.

\begin{table}[H]
\centering
\caption{Experimental configuration matrix.}
\label{tab:experiments}
\footnotesize
\renewcommand{\arraystretch}{1.0}
\begin{tabular}{@{}cllllc@{}}
\toprule
\textbf{Model} & \textbf{Downsample} & \textbf{Blur $\sigma$} & \textbf{Noise} & \textbf{Scale} & \textbf{LR size} \\
\midrule
M1 & Bicubic    & 1.0 & Gaussian $\sigma_n=5$  & $\times$2 & 256$\times$256 \\
M2 & Bicubic    & 1.0 & Gaussian $\sigma_n=5$  & $\times$4 & 128$\times$128 \\
M3 & Bicubic    & 1.0 & Gaussian $\sigma_n=5$  & $\times$8 & 64$\times$64 \\
\midrule
M4 & Bilinear   & 1.5 & Gaussian $\sigma_n=10$ & $\times$2 & 256$\times$256 \\
M5 & Bilinear   & 1.5 & Gaussian $\sigma_n=10$ & $\times$4 & 128$\times$128 \\
M6 & Bilinear   & 1.5 & Gaussian $\sigma_n=10$ & $\times$8 & 64$\times$64 \\
\midrule
M7 & Lanczos    & 2.0 & Poisson                & $\times$2 & 256$\times$256 \\
M8 & Lanczos    & 2.0 & Poisson                & $\times$4 & 128$\times$128 \\
M9 & Lanczos    & 2.0 & Poisson                & $\times$8 & 64$\times$64 \\
\bottomrule
\end{tabular}
\end{table}

Each combination yields one trained model (9 total).
\begin{figure}[H]
\centering
\includegraphics[width=\columnwidth]{degradation_grid.png}
\caption{Visual degradation across all scenarios at $\times$2. HR (512px) $\rightarrow$ LR (256px) through each pipeline.}
\label{fig:degradation}
\end{figure}

\subsection{Generator Architecture}

The generator $G$ is a deep residual network (SRResNet)~\cite{ledig2017}.
It maps LR input $\mathbf{I}_{\text{LR}}$ to SR output $\mathbf{I}_{\text{SR}}$.
The network begins with a 9$\times$9 convolution (64 channels) and PReLU activation.
\textbf{16 residual blocks} follow, each with two 3$\times$3 convolutions, batch normalization, and PReLU, connected via local skip connections.
A post-residual 3$\times$3 convolution with batch normalization adds a global skip connection from the head~\cite{he2016}.
The upsampling module uses $N = \log_2(s)$ PixelShuffle blocks~\cite{shi2016}.
Each block applies a 3$\times$3 convolution (4$\times$ channel expansion), sub-pixel shuffle, and PReLU.
A final 9$\times$9 convolution with $\tanh$ produces the 3-channel output.
PixelShuffle avoids checkerboard artifacts common to transposed convolutions~\cite{shi2016}.

\subsection{Discriminator Architecture}

The discriminator $D$ classifies images as real HR or generated SR.
It follows a VGG-style~\cite{simonyan2015} design with eight convolutional blocks.
Each block contains a 3$\times$3 convolution, batch normalization (except the first), and LeakyReLU ($\alpha=0.2$).
Strided convolutions ($s=2$) halve spatial dimensions every two blocks while doubling channels: 64 $\rightarrow$ 128 $\rightarrow$ 256 $\rightarrow$ 512.
The final feature map is flattened and passed through two fully connected layers (1024 $\rightarrow$ 1) with sigmoid output.

\subsection{Loss Functions}

The generator is optimized with a composite loss:

\begin{equation}
\mathcal{L}_G = \lambda_{\text{content}} \mathcal{L}_{\text{VGG}} + \lambda_{\text{adv}} \mathcal{L}_{\text{adv}} + \lambda_{\text{pixel}} \mathcal{L}_{\text{pixel}}
\label{eq:loss}
\end{equation}

where $\lambda_{\text{content}} = 1.0$, $\lambda_{\text{adv}} = 1 \times 10^{-3}$, and $\lambda_{\text{pixel}} = 1 \times 10^{-2}$~\cite{ledig2017}.
The three terms are:

\textbf{Perceptual loss} $\mathcal{L}_{\text{VGG}}$: MSE between VGG-19~\cite{simonyan2015} ReLU3\_3 activations of $\mathbf{I}_{\text{SR}}$ and $\mathbf{I}_{\text{HR}}$.
This measures semantic similarity, not pixel error~\cite{johnson2016}.

\textbf{Adversarial loss} $\mathcal{L}_{\text{adv}}$: binary cross-entropy where $G$ maximizes $D(G(\mathbf{I}_{\text{LR}}))$ to fool the discriminator~\cite{goodfellow2014}.

\textbf{Pixel loss} $\mathcal{L}_{\text{pixel}}$: L1 distance between $\mathbf{I}_{\text{SR}}$ and $\mathbf{I}_{\text{HR}}$, providing structural stability.

The discriminator is trained independently with binary cross-entropy:

\begin{equation}
\mathcal{L}_D = -\mathbb{E}\left[\log D(\mathbf{I}_{\text{HR}})\right] - \mathbb{E}\left[\log\left(1 - D(G(\mathbf{I}_{\text{LR}}))\right)\right]
\label{eq:disc_loss}
\end{equation}

\subsection{Training Protocol}

Training has two phases.
\textbf{Phase 1 (pre-training)}: the generator trains alone with L1 loss for TBD epochs using Adam~\cite{kingma2014} at $10^{-4}$ learning rate.
This initializes weights in a stable configuration before adversarial training.

\textbf{Phase 2 (full GAN)}: discriminator and generator update alternately for TBD epochs.
The discriminator learns to distinguish real HR from generated SR.
The generator uses the full loss (Eq.~\ref{eq:loss}).
Learning rate halves at the training midpoint~\cite{ledig2017}.

Hardware: NVIDIA GeForce GTX 1650 (4\,GB VRAM), PyTorch 2.1.0, CUDA 11.8.
Batch size: 16.

\subsection{DR Classification}

A ResNet50~\cite{he2016} pre-trained on ImageNet is adapted for ordinal regression.
The classifier head is replaced by a single neuron with MSE loss, discretized to five classes via OptimizedRounder threshold optimization.
Gradual unfreezing proceeds in three phases: head only (10 epochs, LR=$10^{-3}$), layer4 + head (20 epochs, LR=$10^{-4}$), and full fine-tune (30 epochs, LR=$10^{-5}$).
Batch size: 16. Input size: 224$\times$224.
OptimizedRounder~\cite{kingma2014} finds the four class boundaries that maximize QWK on the validation set.
Primary metric: Quadratic Weighted Kappa (QWK).
Secondary metrics: accuracy, precision, recall, sensitivity, specificity.

	
\section{Results}
\label{sec:results}

\subsection{Super-Resolution}

All 9 models (Table~\ref{tab:experiments}) are evaluated on the test set.
Metrics: PSNR, SSIM, LPIPS~\cite{zhang2018}, and MSE.
Table~\ref{tab:metrics} presents the full results.

\begin{table}[H]
\centering
\caption{Quantitative evaluation of all 9 SR models on the test set (Phase 1, L1 loss).}
\label{tab:metrics}
\footnotesize
\renewcommand{\arraystretch}{1.0}
\begin{tabular}{@{}lcccc@{}}
\toprule
\textbf{Model} & \textbf{PSNR (dB)} $\uparrow$ & \textbf{SSIM} $\uparrow$ & \textbf{LPIPS} $\downarrow$ & \textbf{MSE} $\downarrow$ \\
\midrule
M1 & 35.19 & 0.922 & 0.0601 & 3.18e-4 \\
M2 & 30.84 & 0.817 & 0.2064 & 9.12e-4 \\
M3 & 26.93 & 0.666 & 0.4887 & 2.58e-3 \\
\midrule
M4 & \textbf{37.43} & \textbf{0.944} & \textbf{0.0429} & \textbf{1.91e-4} \\
M5 & 32.01 & 0.860 & 0.1379 & 6.96e-4 \\
M6 & 27.49 & 0.729 & 0.4135 & 2.19e-3 \\
\midrule
M7 & 34.79 & 0.918 & 0.0799 & 3.48e-4 \\
M8 & 31.20 & 0.855 & 0.1709 & 8.16e-4 \\
M9 & 27.52 & 0.765 & 0.4100 & 2.09e-3 \\
\bottomrule
\end{tabular}
\end{table}

\textbf{Best model: M4} (Scenario B, $\times$2).
PSNR = 37.43\,dB, SSIM = 0.944, LPIPS = 0.0429, MSE = 1.91e-4.

\subsubsection{Scale Factor Impact}

\begin{itemize}
    \item $\times$\textbf{2} (avg): PSNR 35.80\,dB, SSIM 0.928, LPIPS 0.061 --- \textit{excellent}
    \item $\times$\textbf{4} (avg): PSNR 31.35\,dB, SSIM 0.844, LPIPS 0.172 --- \textit{good}
    \item $\times$\textbf{8} (avg): PSNR 27.31\,dB, SSIM 0.720, LPIPS 0.437 --- \textit{limited}
\end{itemize}

$\times$2 to $\times$4 drops 4.45\,dB. $\times$4 to $\times$8 drops 4.04\,dB.
Higher scale factors are increasingly ill-posed.

\subsubsection{Scenario Comparison}

Scenario B (bilinear) achieves the highest average PSNR (32.31\,dB) and lowest LPIPS (0.198).
Scenario A (bicubic): 30.99\,dB. Scenario C (Lanczos): 31.17\,dB.
Bilinear introduces the least aliasing, producing smoother LR inputs that are easier to reconstruct.
Scenario C maintains competitive SSIM (0.846 avg) despite the most aggressive degradation.

\subsection{Visual Results}

\begin{figure}[H]
\centering
\includegraphics[width=\columnwidth]{sr_comparison.png}
\caption{LR $\rightarrow$ SR $\rightarrow$ HR for best ($\times$2, top) and worst ($\times$8, bottom) models. Zoom patches show fine detail recovery.}
\label{fig:sr_results}
\end{figure}

At $\times$2, SR outputs closely approximate HR reference.
Vessel structures and optic disc boundaries are visibly recovered.
At $\times$8, fine vessels are not fully restored, but overall fundus structure and lesion contrast are preserved.

\subsection{DR Classification}

ResNet50 ordinal regression achieves 73.56\% accuracy and QWK = 0.8764 on the test set (537 images).
This improves over the VGG16 baseline (64.43\%) by +9.13\%.

\begin{table}[H]
\centering
\caption{ResNet50 ordinal regression per-class results on the test set.}
\label{tab:classification}
\footnotesize
\renewcommand{\arraystretch}{1.0}
\begin{tabular}{@{}lcccc@{}}
\toprule
\textbf{Class} & \textbf{Support} & \textbf{Precision} & \textbf{Recall} & \textbf{F1-score} \\
\midrule
No DR (0)        & 271 & 0.960 & 0.978 & 0.969 \\
Mild (1)         &  53 & 0.518 & 0.547 & 0.532 \\
Moderate (2)     & 143 & 0.728 & 0.525 & 0.610 \\
Severe (3)       &  27 & 0.232 & 0.815 & 0.361 \\
Proliferative (4)&  43 & 0.571 & 0.093 & 0.160 \\
\midrule
Weighted avg     & 537 & 0.787 & 0.736 & 0.735 \\
\bottomrule
\end{tabular}
\end{table}

Binary DR detection (classes 1--4 vs.~0) reaches 96.83\% accuracy, 95.86\% sensitivity, and 97.79\% specificity.
Class 0 (no DR) is reliably classified (F1 = 0.969).
Minority classes (mild, severe, proliferative) show lower F1-scores due to limited training samples (5.1\%--9.8\%).

	
\section{Discussion}
\label{sec:discussion}

\subsection{Key Findings}

$\times$2 models exceed 35\,dB PSNR across all scenarios.
This indicates moderate SR (256$\times$256 $\rightarrow$ 512$\times$512) is well-posed for retinal fundus images.
Regular vascular structure and optic disc geometry provide strong spatial priors.

Scenario B (bilinear) consistently outperforms A and C.
Bilinear's 4-pixel averaging kernel minimizes aliasing in the LR domain.
The network learns a simpler mapping function.
This suggests smoother downsampling is preferable for reconstruction fidelity over edge-preserving methods in retinal SR.

LPIPS reveals a complementary pattern.
Scenario C achieves competitive perceptual quality at $\times$4 and $\times$8 despite lower PSNR.
Deep feature similarity is less sensitive to the degradation model than to scale factor~\cite{zhang2018}.

\subsection{Limitations}

Full adversarial SRGAN training (generator + discriminator + perceptual loss) could not complete on the available GPU (4\,GB VRAM).
Reported results correspond to a generator trained with L1 loss only.
Full SRGAN training is ongoing and will be reported upon completion.

ResNet50 ordinal regression achieves 73.56\% accuracy and QWK = 0.8764 on five-class DR grading.
This outperforms VGG16 by +9.13\%.
Ordinal regression captures the ordered nature of DR severity better than categorical cross-entropy~\cite{he2016}.
Binary DR detection is excellent (96.83\% accuracy, 95.86\% sensitivity).
Fine-grained severity discrimination remains challenging for minority classes.
Super-resolution alone cannot resolve all grading ambiguities, especially for early-stage microaneurysms.

This study does not include clinical validation with human experts.
PSNR, SSIM, and LPIPS are objective but do not measure clinical utility.
A reader study with ophthalmologists is needed to validate diagnostic impact.

\subsection{Comparison with Prior Work}

The 9-scenario framework extends prior retinal SR studies by varying both degradation and scale within a unified design.
Most prior work evaluates a single pipeline, preventing cross-study comparison~\cite{zhang2020}.
Our matrix provides a reproducible benchmark for future retinal SR research.

ResNet50 results align with Gulshan et al.~\cite{gulshan2016}: deep models detect referable DR well (96.83\% binary accuracy) but struggle with severity grading.
Ordinal regression improves five-class accuracy over VGG16 from 64.43\% to 73.56\%.

	
\section{Conclusions}
\label{sec:conclusions}

This paper presented a systematic SRGAN framework for retinal fundus super-resolution across 9 degradation scenarios.

Key findings:
\begin{itemize}
    \item $\times$2 SR achieves PSNR $>$ 35\,dB and SSIM $>$ 0.92, effectively compensating for moderate resolution loss.
    \item Bilinear degradation (Scenario B) yields the highest metrics across all scales.
    \item Scale factor dominates reconstruction difficulty: $\times$8 models reach only $\sim$27\,dB PSNR.
    \item ResNet50 ordinal regression achieves 73.56\% accuracy and QWK = 0.8764 on five-class DR grading.
    \item Binary DR detection reaches 96.83\% accuracy, 95.86\% sensitivity, and 97.79\% specificity.
\end{itemize}

Future work includes:
\begin{itemize}
    \item Completion of full adversarial SRGAN training with the complete loss function (Eq.~\ref{eq:loss}).
    \item Attention mechanisms for fine-detail recovery at high scale factors.
    \item Class-balanced training and data augmentation for minority DR grades.
    \item Lightweight architectures (e.g., MobileNet) for deployment on portable fundus cameras.
    \item Clinical reader studies to validate diagnostic utility of SR-enhanced images.
\end{itemize}

Code, models, and configuration files: \url{https://github.com/JuanGamez99/ProyectoAvanzado2}.


%%%% Create the references below (DO NOT CHANGE) 
%%%% (MORE INFORMATION IN README.txt) %%%%

%%%%%%%%%%%%%%%%%%%%%%%%%%%%%%% ATTENTION %%%%%%%%%%%%%%%%%%%%%%%%%%%%%%%%
%%%% The file extra considerations (extra_considerations.tex) is not part of the template and must be removed from the final work. To do so, comment the line below %%%%
\input{extra_considerations.tex} 
\newpage
\renewcommand{\refname}{REFERENCES}

\begin{thebibliography}{99}

\bibitem{1}
D.~A. Chiran Taimal, ``Optimización del monitoreo ambiental mediante aplicaciones digitales en reforestaciones y restauraciones ecológicas en Cumbal, Colombia,'' 2024.


\bibitem{2}
C.~A. Pedraza Peñaloza, \textit{Novel methods and technologies to improve monitoring and understanding of land use land cover dynamics based on satellite Earth Observation: The case of forest monitoring in Colombia}, 2023.

\bibitem{3}
J.~A. Achicanoy, R. Rojas-Robles, and J.~E. Sánchez, ``Land cover vegetation analysis and projection through remote sensing and Geographic Information Systems in the Suba District, Bogotá-Colombia,'' \textit{Gestión y Ambiente}, vol.~21, no.~1, p.~41, 2018.

\bibitem{4}
M.~O. Chávez Can, \textit{Uso de imágenes de satélite de alta resolución espacial en la detección de cultivos agrícolas en el municipio de Patzicía, Chimaltenango}, Ph.D. dissertation, Univ. de San Carlos de Guatemala, Guatemala City, Guatemala, 2021.


\bibitem{2012}
S. Khorram, F. H. Koch, C. F. Van der Wiele, and S. A. Nelson, 
\textit{Remote Sensing}. Springer Science \& Business Media, 2012.

\bibitem{5}
A.~Vijayakumar and S.~Vairavasundaram, 
``YOLO-based object detection models: A review and its applications,'' 
\textit{Multimedia Tools and Applications}, vol.~83, pp.~83535--83574, 2024. [Online]. Available: \url{https://doi.org/10.1007/s11042-024-18872-y}

\bibitem{6}
European Space Agency (ESA) y European Commission,
``Copernicus Sentinel‑2: a multispectral high‑resolution land monitoring mission,'' 
\textit{Copernicus Programme}, 2024. [En línea]. Disponible en: \url{https://sentinels.copernicus.eu/copernicus/sentinel-2}

\bibitem{7}
J.~D.~Lencinas y A.~Siebert, 
``Relevamiento de bosques con información satelital: Resolución espacial y escala,'' 
\textit{Quebracho - Revista de Ciencias Forestales}, vol.~17, no.~1-2, pp.~101--105, 2009.


\bibitem{8}
Instituto Geográfico Agustín Codazzi (IGAC), 
``Observatorio de la Tierra y el Territorio,'' 
\textit{IGAC}, 2024. [En línea]. Disponible en: \url{https://igacplanet.azurewebsites.net/}



\bibitem{9}
SIATA - Sistema de Alerta Temprana de Medellín y el Valle de Aburrá, 
``Nosotros – ¿Qué es el SIATA?,'' 
\textit{SIATA}, 2025. [En línea]. Disponible en: \url{https://siata.gov.co/sitio_web/index.php/nosotros}


\bibitem{10}
P.~S.~P. Reyes, A.~S.~V. Maholy, R.~P. Cabrera, R.~O.~S.~I. Verdezoto, B.~A.~C. Macías, and K.~M. Veliz, 
``Determinación de la pérdida de cobertura vegetal en el Cantón Manta (Periodo 2000–2022) / Determination of the loss of plant coverage in the Manta Canton (Period 2000–2022),'' 
\bibitem{11}
M.~Reyes and I.~Lizarazo, 
``Convolutional neural network-based semantic segmentation model for land cover classification in páramo ecosystems,'' 
\textit{Revista de Teledetección}, no.~65, 2025. [Online]. Available: \url{https://doi.org/10.4995/raet.2025.21858}



\textit{2024}.
\bibitem{12}
S. Razakarivony and F. Jurie, 
``Vehicle detection in aerial imagery: A small target detection benchmark,'' 
\textit{Journal of Visual Communication and Image Representation}, vol. 34, pp. 187--203, 2015.
\bibitem{13}
P. Ebel, A. Meraner, M. Schmitt, and X. X. Zhu, 
``Multisensor data fusion for cloud removal in global and all-season Sentinel-2 imagery,'' 
\textit{IEEE Transactions on Geoscience and Remote Sensing}, 2020.

\end{document} % end of document %%%%(DO NOT CHANGE)%%%%
